\documentclass[11pt,a4paper]{article}

\usepackage[utf8]{inputenc}
\usepackage[T1]{fontenc}
\usepackage[margin=1in]{geometry}
\usepackage{amsmath,amssymb,amsthm}
\usepackage{booktabs}
\usepackage{enumitem}
\usepackage{longtable,array,url,placeins}
\usepackage[hidelinks]{hyperref}
\usepackage[protrusion=true,expansion=false]{microtype}

\theoremstyle{plain}
\newtheorem{theorem}{Theorem}
\newtheorem{proposition}[theorem]{Proposition}
\theoremstyle{definition}
\newtheorem{definition}[theorem]{Definition}

\newcommand{\R}{\mathbb{R}}
\newcommand{\E}{\mathbb{E}}
\newcommand{\Var}{\operatorname{Var}}
\newcommand{\simplex}{\Delta^{m-1}}
\newcommand{\pipre}{\pi^{\mathrm{pre}}}
\newcommand{\pieq}{\pi^{\mathrm{eq}}}
\newcommand{\pir}{\pi^{r}}
\newcommand{\pihat}{\widehat{\pi}}
\newcommand{\epsincons}{\varepsilon_{\mathrm{incons}}}
\newcommand{\epssurr}{\varepsilon_{\mathrm{surr}}}
\newcommand{\epslearn}{\varepsilon_{\mathrm{learn}}}

\usepackage{graphicx}

\newcolumntype{L}[1]{>{\raggedright\arraybackslash}p{#1}}
\setlength{\emergencystretch}{2em}
\renewcommand{\topfraction}{0.9}
\renewcommand{\bottomfraction}{0.85}
\renewcommand{\textfraction}{0.1}
\renewcommand{\floatpagefraction}{0.7}
\title{\vspace{-1.5em}What Additive Risk-Adjusted Rewards Actually Optimise\\
in Reinforcement Learning for Asset Allocation}
\author{Matin Tavakoli}
\date{Research Proposal}
\begin{document}
\maketitle
\vspace{-1.5em}


\begin{abstract}
Financial reinforcement learning translates an investment objective into a reward that an agent can optimise. This translation matters: a policy can maximise its training reward while pursuing a different risk--return trade-off from the one the investor intended. This project studies that difference in a portfolio model where both the desired policy and the reward-optimal policy can be calculated. Taking terminal mean--variance equilibrium as the target, we derive the policy mismatch that remains after optimally calibrating a quadratic wealth-increment reward. We then construct a reward based on excess gains measured in terminal units that recovers the target under independent returns and unconstrained dollar holdings. Analytical results and controlled experiments connect reward choice to portfolio behaviour, distinguish reward-design error from learning error, and provide a practical basis for calibrating financial RL rewards.
\end{abstract}

\section{Motivation and research question}
A portfolio manager may care about wealth at a future date, while an RL agent receives rewards after each trade. The reward must connect these two perspectives. Penalising squared changes in wealth is an appealing choice: it is simple to implement and encourages a balance between gains and risk. Yet the policy induced by this reward need not coincide with the policy implied by a terminal financial objective. Better optimisation of the training reward does not resolve this difference.

The central question is therefore: \emph{when does an additive risk-adjusted reward induce the intended portfolio policy, and how should it be changed when it does not?} We address this question using terminal mean--variance preferences. These preferences distinguish the policy chosen at the initial date from the equilibrium policy that successive decision makers are willing to follow. The project takes the latter as its design target and retains the initial-date optimum as a separate performance reference.

Our main result quantifies a structural obstacle. With nonzero interest, one scalar penalty on squared wealth increments generally cannot recover the equilibrium policy. For every fixed horizon, we derive the exact leading coefficient of the mismatch that remains after choosing the best penalty. A reward expressed in terminal units removes this obstacle in the independent-return model. Figure~\ref{fig:alignment} illustrates the resulting change in portfolio decisions.

\begin{figure}[htbp]
\centering\includegraphics[width=\linewidth]{../figures/01_policy_alignment.pdf}
\caption{Portfolio policies under alternative rewards, with three risky assets, $T=4$, $\varphi=1$ and risk-free gross return $s=1.02$. Scalar calibration reduces the distance to equilibrium; the terminal-unit reward recovers it. Policy alignment and initial-date mean--variance value measure different aspects of performance.}
\label{fig:alignment}
\end{figure}

\section{Related research}
Financial RL uses several approaches to reward design. Kolm and Ritter study an approximate quadratic increment reward for hedging and explicitly discuss the approximation involved \cite{ref1}. Jiang, Xu and Liang optimise log portfolio growth \cite{ref2}, while Yang et al. train trading agents on changes in portfolio value and use Sharpe ratios for ensemble selection \cite{ref3}. Moody and Saffell derive differential Sharpe rewards for direct reinforcement learning \cite{ref4}. These choices express different preferences, making the relationship between training reward and evaluation objective an important part of interpreting a learned strategy.

Risk-sensitive RL also provides explicit reward transformations. Bisi et al. analyse per-step reward volatility, and Zhang et al. develop mean--variance policy iteration \cite{ref5,ref6}. In our zero-interest independent-return benchmark, an adaptation of the latter has the same fixed point as the scalar calibration derived here. The contribution therefore centres on the quantitative mismatch of the raw reward family and the conditions under which that mismatch can be removed.

The financial target has a well-developed foundation. Li and Ng derive the multiperiod precommitted mean--variance solution, while Basak and Chabakauri study time-consistent asset allocation \cite{ref9,ref10}. Wang and Zhou formulate RL directly for continuous-time mean--variance selection, and Dai, Dong and Jia learn equilibrium strategies for a log-return mean--variance objective \cite{ref7,ref8}. Building on these approaches, we examine the intermediate link between an investor's chosen target and a convenient additive training reward. We hold the market and feasible actions fixed, solve the reward-induced problem, and measure its policy distance from the target before introducing learning.

\section{From the investment objective to the reward}
Consider a finite investment horizon $T$. Wealth evolves according to
\begin{equation}
 X_{t+1}=s_tX_t+P_t^{\top}u_t,\qquad
 J(\pi)=\E[X_T]-\varphi\Var(X_T),\quad\varphi>0,
\end{equation}
where $s_t>0$ is the deterministic risk-free gross return, $P_t$ is the vector of excess returns, and $u_t$ denotes risky dollar holdings. The analytical benchmark assumes independent returns with known means $m_t$ and positive-definite covariances $\Sigma_t$, and unconstrained dollar positions. Finite second moments suffice for the exact policy calculations.

Define
\begin{equation}
 M_t=\Sigma_t+m_tm_t^{\top},\quad d_t=M_t^{-1}m_t,\quad
 \nu_t=m_t^{\top}d_t<1,\quad H_t=\prod_{j=t+1}^{T-1}s_j.
\end{equation}
The factor $H_t$ converts a gain at the next date into terminal wealth. The equilibrium risky allocation is
\begin{equation}
 u_t^{\mathrm{eq}}=\frac{\Sigma_t^{-1}m_t}{2\varphi H_t}.
\end{equation}
This policy provides a precise target against which reward choices can be assessed.

The proposed procedure starts with this target, solves the control problem induced by the candidate reward, and selects reward parameters to minimise policy distance. When scalar calibration leaves a structural mismatch, we redesign the reward. For the independent-return benchmark, the resulting rule is
\begin{equation}
 Y_t=H_tP_t^{\top}u_t,\qquad
 r_t^{\mathrm{design}}=Y_t-\varphi(1-\nu_t)Y_t^2.
\end{equation}
Its conditional expected reward is strictly concave in the action and maximised at $u_t^{\mathrm{eq}}$. Independence and wealth-independent feasible dollar holdings make these pointwise maximisers optimal for the additive criterion. The construction combines terminal-unit scaling with a correction for the difference between a second moment and a variance. When moments change over time, the correction changes with them.

\subsection{Why scalar calibration is not enough}
For iid returns, write $d=M^{-1}m$, $\nu=m^{\top}d$ and $b=[2\varphi(1-\nu)]^{-1}$. Consider the raw reward
\[
 r_t=\Delta X_t-\kappa(\Delta X_t)^2,\qquad\kappa>0.
\]
Let $D$ be the square root of the sum of expected squared action differences from equilibrium, evaluated along the reward-optimal policy. For nonzero $m$, fixed $T\ge2$ and $s$ near one, we obtain
\begin{equation}
\begin{aligned}
 \min_{\kappa>0}D(\pi^\kappa,\pieq)&=c_T|s-1|+O((s-1)^2),\\
 c_T&=\|d\|b\sqrt{\frac{T(T^2-1)}{3}
 +\frac{\nu(1-\nu)T(T-1)}{2}}.
\end{aligned}
\end{equation}
At zero interest, $\kappa=\varphi(1-\nu)$ gives exact recovery. Away from zero interest, the final reward-optimal allocation has wealth slope $-d(s-1)$, whereas equilibrium dollar holdings are independent of wealth. Changing the scalar penalty cannot eliminate this difference across reachable wealth states.

The coefficient makes the size of this obstruction explicit. Its leading growth with horizon is $T^{3/2}$, highlighting the importance of reward choice over longer investment periods. The expansion is local in interest for each fixed horizon. Figure~\ref{fig:horizon} compares it with exact calibration.

\begin{figure}[htbp]
\centering\includegraphics[width=\linewidth]{../figures/02_horizon_theorem.pdf}
\caption{How the investment horizon affects the best achievable scalar calibration. Solid curves show exact calculations and dashed curves show the local analytical prediction. The coefficient grows as $T^{3/2}$; the expansion is taken at a fixed horizon.}
\label{fig:horizon}
\end{figure}

\subsection{Adapting to changing market conditions}
Even at zero interest, a common penalty can be insufficient when return moments change. With $a_t=\|d_t\|^2$, $b_t=[2\varphi(1-\nu_t)]^{-1}$ and $\bar b=\sum_ta_tb_t/\sum_ta_t$, the minimum squared mismatch is
\begin{equation}
 \min_{\kappa>0}D^2=\sum_ta_t(b_t-\bar b)^2.
\end{equation}
This expression identifies when a date-specific calibration is needed: the optimal penalty must vary across dates with nonzero weight. It also provides a direct measure of the cost of retaining one common coefficient.

\section{Evidence from controlled portfolio experiments}
The analytical benchmark makes it possible to compare rewards without confounding reward choice with training quality. In the three-asset example of Figure~\ref{fig:alignment}, the raw penalty produces policy distance $0.797127$. Zero-interest scalar calibration reduces it to $0.195347$, and the best scalar choice reduces it to $0.186435$. The terminal-unit reward reduces it to numerical zero. Across 120 combinations of horizon, risk aversion and interest, the largest discrepancy between the redesigned and equilibrium policies is $2.35\times10^{-14}$.

Financial performance is evaluated alongside alignment. In the same example, the raw reward has terminal mean--variance value $1.316580$, compared with $1.323644$ for equilibrium, yet its excess Sharpe ratio is higher: $0.987285$ versus $0.982266$. The best scalar reward has value $1.326093$, slightly above equilibrium despite its nonzero policy distance. These comparisons show why the target must be specified in advance: matching equilibrium, maximising initial-date value and maximising Sharpe are distinct goals.

\subsection{Comparing rewards used in financial RL}
A second benchmark uses a common long-only action set to compare value-change, log-growth, quadratic and differential Sharpe rewards. It has three decision periods, initial wealth one, risk-free gross return $1.02$, and risky excess returns $-0.18$ or $0.30$ with equal probability. Each reward is solved by finite-tree dynamic programming on the same 21-point weight grid; equilibrium is computed on that grid as well. Differential Sharpe uses its full moving-moment state.

\begin{figure}[htbp]
\centering\includegraphics[width=\linewidth]{../figures/04_published_rewards.pdf}
\caption{Reward formulations adapted to a common long-only market. The comparison concerns induced portfolio policies rather than reproductions of the original papers' algorithms or backtests. The quadratic coefficient $0.05$ is included as a wealth-unit-dependent sensitivity example from Kolm--Ritter.}
\label{fig:rewards}
\end{figure}

The terminal-unit design reduces policy distance from $0.058485$ for the raw reward to $0.041366$ in this constrained example. Its terminal mean--variance value is slightly lower ($1.107729$ versus $1.107849$). Wealth-dependent position constraints change the control problem, so the unconstrained recovery result becomes a starting point for constrained reward design. Log growth and value-change rewards provide useful comparisons between investment preferences, rather than alternative implementations of the same objective.

\subsection{Returns with observable dependence}
The method also identifies what must change when future investment opportunities depend on the current return. In a scalar market with an observed regime $z$, the equilibrium allocation includes a continuation hedge:
\begin{equation}
 u^{\mathrm{eq}}(t,z)=\frac{m_z}{2\varphi H_t\Var(P\mid z)}
 -\frac{\operatorname{Cov}(P,a_{t+1}(z_{t+1})\mid z)}{H_t\Var(P\mid z)},
\end{equation}
where $a_{t+1}$ is expected terminal continuation gain. Current conditional moments alone omit the second term. Using the solved continuation, adding the linear correction
$-2\varphi H_tu\operatorname{Cov}(P,a_{t+1}\mid z)$ to the reward recovers the equilibrium in this model. This implements the intertemporal hedging principle within the reward construction \cite{ref10}.

\begin{figure}[htbp]
\centering\includegraphics[width=\linewidth]{../figures/05_dependent_returns.pdf}
\caption{Reward calibration with dependent returns. Exact recursion over the joint return and next-regime distribution identifies the continuation hedge. At persistence $0.8$, policy distance is $0.365102$ for the raw reward and $0.317644$ for local calibration; the hedge-corrected reward recovers equilibrium to numerical precision.}
\label{fig:dependent}
\end{figure}

\section{Reward design and learning}
Knowing the reward-optimal policy provides a reference for evaluating a trained agent. We distinguish the effect of time consistency, the policy induced by the reward, and the remaining learning difference through
\begin{equation}
\begin{aligned}
 J(\pipre)-J(\pihat)
 &= [J(\pipre)-J(\pieq)]\\
 &\quad+[J(\pieq)-J(\pir)]
 +[J(\pir)-J(\pihat)].
\end{aligned}
\end{equation}
The latter two terms are signed financial-value differences. We also measure surrogate regret, the nonnegative difference between the expected training reward at its exact optimum and at the learned policy. Together, these measures distinguish an inappropriate reward from incomplete optimisation of an appropriate one.

\begin{figure}[htbp]
\centering\includegraphics[width=\linewidth]{../figures/06_learning_error.pdf}
\caption{Initial learning experiments with five paired seeds per reward, a common affine policy family, 4,000 iterations and batches of 512. The exact reward optima provide references for interpreting the learned policies and separating optimisation error from reward mismatch.}
\label{fig:learning}
\end{figure}

In this initial study, mean policy distance falls from $0.8638$ under the raw reward to $0.8079$ under the design. Mean terminal value is $1.305628$ and $1.300671$, respectively. The result illustrates why an exact alignment guarantee must be evaluated together with learning behaviour: it determines the destination of optimisation, while the algorithm and training budget determine how closely that destination is reached.

\section{Research programme and significance}
The next stage will examine how this approach transfers from exact models to implementable portfolio systems. First, constrained allocations and transaction costs will be incorporated using a shared feasible benchmark, so that changes in reward can be distinguished from changes in the trading problem. Second, the continuation hedge will be estimated from observed data rather than supplied by the model. Third, matched-budget policy-gradient, PPO and DDPG experiments will assess how reward choice interacts with learning, using separate validation data for tuning and repeated runs for uncertainty estimates.

Robustness will be assessed across risk aversion, horizon, estimated moments and return dependence. Existing moment-estimation and stress experiments provide a starting point for this stage: they show that serial dependence can change the value ranking of policies calibrated in an iid market. Historical evaluation will therefore report policy alignment where a reference is available, terminal mean--variance performance, turnover and learning stability alongside returns and Sharpe ratios. Differential Sharpe rewards will be studied with their moving-moment state retained throughout the control problem.

The practical contribution is a way to make reward choice an explicit part of financial model design. A researcher can calculate what a reward asks the agent to do, measure how far that policy lies from the intended objective, and determine whether a better penalty is sufficient or a different reward is required. The analytical mismatch result and the recovery construction make those decisions testable before committing to a large training or backtesting exercise.
\FloatBarrier

\clearpage
\begin{thebibliography}{99}
\bibitem{ref1} Petter N. Kolm; Gordon Ritter. Dynamic Replication and Hedging: A Reinforcement Learning Approach. The Journal of Financial Data Science 1(1), 159--171 (2019). \url{https://doi.org/10.3905/jfds.2019.1.1.159}

\bibitem{ref2} Zhengyao Jiang; Dixing Xu; Jinjun Liang. A Deep Reinforcement Learning Framework for the Financial Portfolio Management Problem. arXiv preprint 1706.10059 (2017). \url{https://arxiv.org/abs/1706.10059}

\bibitem{ref3} Hongyang Yang; Xiao-Yang Liu; Shan Zhong; Anwar Walid. Deep Reinforcement Learning for Automated Stock Trading: An Ensemble Strategy. Proceedings of the First ACM International Conference on AI in Finance (ICAIF), 8 pages (2020). \url{https://doi.org/10.1145/3383455.3422540}

\bibitem{ref4} John Moody; Matthew Saffell. Learning to Trade via Direct Reinforcement. IEEE Transactions on Neural Networks 12(4), 875--889 (2001). \url{https://doi.org/10.1109/72.935097}

\bibitem{ref5} Lorenzo Bisi; Luca Sabbioni; Edoardo Vittori; Matteo Papini; Marcello Restelli. Risk-Averse Trust Region Optimization for Reward-Volatility Reduction. Proceedings of IJCAI-20, 4583--4589 (2020). \url{https://doi.org/10.24963/ijcai.2020/632}

\bibitem{ref6} Shangtong Zhang; Bo Liu; Shimon Whiteson. Mean-Variance Policy Iteration for Risk-Averse Reinforcement Learning. Proceedings of AAAI 35(12), 10905--10913 (2021). \url{https://doi.org/10.1609/aaai.v35i12.17302}

\bibitem{ref7} Haoran Wang; Xun Yu Zhou. Continuous-Time Mean-Variance Portfolio Selection: A Reinforcement Learning Framework. Mathematical Finance 30(4), 1273--1308 (2020). \url{https://doi.org/10.1111/mafi.12281}

\bibitem{ref8} Min Dai; Yuchao Dong; Yanwei Jia. Learning Equilibrium Mean-Variance Strategy. Mathematical Finance 33(4), 1166--1212 (2023). \url{https://doi.org/10.1111/mafi.12402}

\bibitem{ref9} Duan Li; Wan-Lung Ng. Optimal Dynamic Portfolio Selection: Multiperiod Mean-Variance Formulation. Mathematical Finance 10(3), 387--406 (2000). \url{https://doi.org/10.1111/1467-9965.00100}

\bibitem{ref10} Suleyman Basak; Georgy Chabakauri. Dynamic Mean-Variance Asset Allocation. The Review of Financial Studies 23(8), 2970--3016 (2010). \url{https://doi.org/10.1093/rfs/hhq028}

\end{thebibliography}

\end{document}
